% =====================================================================
\chapter{Decision Tree (ID3)}
% =====================================================================

\labinfo{12\_Decision\_Tree\_ID3}{id3.py}{UCI Mushroom: 8,124 records, 22 categorical attributes}
{Binary classification (edible / poisonous)}{Test accuracy 100\%, macro F1 1.000}

\section{Dataset}
The \textbf{Mushroom} dataset (UCI) contains 8,124 mushrooms described by 22 categorical attributes.
The class is edible or poisonous. The attributes include cap, gill, stalk, ring, spore-print, population
and habitat properties. Missing categorical values are kept as their own value, so no numerical encoding
or imputation is needed.

Since every attribute is categorical, the data suits ID3 perfectly; no discretisation is needed. I used a
stratified 70/30 split (5,687 training and 2,437 test records).

\section{Code and experiments}
\begin{itemize}
\item \code{entropy} and \code{information\_gain} are direct translations of the two formulas;
      \code{information\_gain} uses \code{groupby} to form the subsets $S_v$.
\item \code{id3} builds the tree recursively as nested dictionaries. An inner node stores the attribute,
      its gain, the majority class and one child per value; a leaf stores a class and the number of
      training examples that reached it. An optional \code{max\_depth} turns it into a depth-limited tree.
\item \code{predict\_one} walks down the tree; if it meets a value it has no branch for, it returns the
      majority class of the current node.
\item \code{tree\_size} and \code{print\_tree} report the size of the tree and print its top levels.
\end{itemize}
\noindent Experiments: (1)~the information gain of every attribute at the root; (2)~the full ID3 tree
on the test set; (3)~trees with maximum depth 1 to 6.

\section{Results}

The entropy of the training set is $H(S)=0.9991$ bits. Table~\ref{tab:id3-gain} shows the gain of each
attribute at the root.

\begin{table}[H]
\centering\small
\caption{Information gain at the root (bits).}
\label{tab:id3-gain}
\begin{tabular}{@{}l c c c c c c@{}}
\toprule
Attribute & \textbf{odor} & spore\_print\_color & gill\_color & ring\_type & stalk\_surface\_above\_ring & stalk\_surface\_below\_ring\\
\midrule
Gain & \textbf{0.9122} & 0.4901 & 0.4206 & 0.3241 & 0.2923 & 0.2883\\
\bottomrule
\end{tabular}
\end{table}

Odor gives the largest gain (0.9122 bits), producing nearly pure edible and poisonous branches. Spore
print colour and gill colour provide the next strongest splits.

The full tree has 5 decision nodes and 33 leaves, with maximum depth 4. It classifies the training set
and test set perfectly in this split.

\begin{table}[H]
\centering\small
\caption{Test-set results of the full tree (2,437 mushrooms).}
\label{tab:id3-test}
\begin{tabular}{@{}l S[table-format=1.3] S[table-format=1.3] S[table-format=1.3] S[table-format=3.0]@{}}
\toprule
\textbf{Class} & {\textbf{Precision}} & {\textbf{Recall}} & {\textbf{F1}} & {\textbf{Support}}\\
\midrule
edible    & 1.000 & 1.000 & 1.000 & 1262\\
poisonous & 1.000 & 1.000 & 1.000 & 1175\\
\midrule
Macro average    & 1.000 & 1.000 & 1.000 & 2437\\
\midrule
\multicolumn{5}{@{}l}{Accuracy: \textbf{1.000} (all 2,437 test mushrooms correct)}\\
\bottomrule
\end{tabular}
\end{table}

The tree is right for 92.3\% of the unseen cars. The two big classes are handled well (F1 0.97 and 0.87),
while the two small classes, good and vgood, are harder (F1 around 0.7). They have only about 45 training
examples each, spread over many leaves, so many of their leaves rest on just one or two cars. Most
mistakes are ``near misses'' between adjacent grades, which is reassuring.

\begin{table}[H]
\centering\small
\caption{Limiting the depth of the tree.}
\label{tab:id3-depth}
\begin{tabular}{@{}l c c c c c c@{}}
\toprule
Maximum depth & 1 & 2 & 3 & 4 & 5 & 6 (full)\\
\midrule
Leaves & 9 & 17 & 23 & 33 & 33 & 33\\
Training accuracy & 0.987 & 0.994 & 0.997 & 1.000 & 1.000 & 1.000\\
Test accuracy & 0.982 & 0.994 & 0.998 & \textbf{1.000} & \textbf{1.000} & \textbf{1.000}\\
\bottomrule
\end{tabular}
\end{table}

The depth study shows the classic overfitting picture only mildly. Because this dataset was generated by
exact rules, deep trees are needed and the test accuracy keeps rising up to depth 5. The last level,
however, adds 87 leaves and takes the training accuracy from 97.5\% to 100\% without improving the test
accuracy at all. A depth-5 tree is therefore just as good and 40\% smaller.

\begin{figure}[H]
\centering
\includegraphics[width=0.48\linewidth]{\figpath{12\_Decision\_Tree\_ID3}{root_gains.pdf}}\hfill
\includegraphics[width=0.48\linewidth]{\figpath{12\_Decision\_Tree\_ID3}{depth_vs_accuracy.pdf}}
\caption{Root information gains and depth study for the Mushroom ID3 tree.}
\end{figure}

\section{Source code}
\lstinputlisting[caption={\texttt{12\_Decision\_Tree\_ID3/id3.py}}]{12_Decision_Tree_ID3/id3.py}
\clearpage
